\newpage
\chapter{TINJAUAN PUSTAKA} \label{Bab II}

\section{Tinjauan Pustaka} \label{II.Tinjauan}
Evaluasi kinerja terhadap sistem antrean pesan, mekanisme penjadwalan tugas, serta perbandingan antara penyimpanan berbasis memori dan basis data relasional telah dikaji dari berbagai sudut pandang pada sejumlah penelitian terdahulu. Tinjauan pustaka ini menelaah lima penelitian terkini untuk memetakan metode pengujian yang digunakan, temuan empiris yang dihasilkan, serta batasan permasalahan yang belum terjawab sehingga membentuk celah penelitian pada Tugas Akhir ini. \par

\begin{enumerate}[noitemsep]
	\item \textbf{Benchmarking Message Queues} oleh Rojina Maharjan, Md Showkat Hossain Chy, Muhammad Ashfakur Arju, dan Tomas Cerny (2023) \cite{maharjan2023benchmarking}. Penelitian ini berangkat dari permasalahan sulitnya menentukan sistem antrean pesan yang tepat di tengah beragamnya arsitektur \textit{message broker} modern. Penulis melakukan eksperimen komparatif menggunakan OpenMessaging Benchmark Framework untuk menguji empat teknologi antrean populer, yaitu Redis, RabbitMQ, ActiveMQ Artemis, dan Apache Kafka, di bawah variasi ukuran pesan serta jumlah \textit{producer} dan \textit{consumer}. Hasil pengujian menunjukkan bahwa Redis mampu memberikan latensi pengiriman paling rendah karena seluruh operasi antreannya berlangsung di dalam memori, sedangkan Apache Kafka unggul jauh pada capaian \textit{throughput} untuk volume pesan masif. Kendati demikian, pengujian Maharjan dkk. hanya mengevaluasi pengiriman pesan instan (\textit{immediate queuing}) dengan berfokus pada metrik \textit{throughput} dan latensi dari sisi luar, tanpa menguji penjadwalan tugas tertunda berbasis waktu (\textit{delayed task scheduling}), ketepatan terhadap target waktu eksekusi (\textit{temporal jitter}), maupun profil konsumsi sumber daya perangkat keras saat terjadi lonjakan beban.
	\item \textbf{DBOS: A DBMS-Oriented Operating System} oleh Athinagoras Skiadopoulos, Qian Li, Peter Kraft, Kostis Kaffes, Michael Cafarella, Michael Stonebraker, Matei Zaharia, dkk. (2022) \cite{skiadopoulos2022dbos}. Penelitian ini mengangkat masalah tingginya kompleksitas pengelolaan status dan pemulihan kegagalan apabila penjadwalan tugas serta komunikasi antarproses dipisahkan ke banyak komponen eksternal. Penulis mengusulkan pendekatan penyimpanan status penjadwalan (\textit{task scheduling}) dan antrean pesan langsung di dalam tabel basis data relasional transaksional menggunakan pengindeksan B-Tree dan transaksi ACID. Hasil evaluasi membuktikan bahwa manajemen jadwal berbasis tabel relasional mampu mencapai latensi penjadwalan yang kompetitif sekaligus menyederhanakan pemulihan pascakegagalan (\textit{fault recovery}) karena status eksekusi tersimpan secara konsisten. Penelitian ini menjadi landasan penting bagi pendekatan \textit{Database-Backed Worker}, namun evaluasi yang dilakukan belum membandingkan secara langsung beban CPU, memori, dan \textit{jitter} dari kueri \textit{polling} non-blokir PostgreSQL (\texttt{FOR UPDATE SKIP LOCKED}) terhadap mesin penjadwal \textit{in-memory} murni maupun antrean berbasis Redis pada spesifikasi server tunggal yang setara.
	\item \textbf{Comparative Evaluation of Java Virtual Machine-Based Message Queue Services: A Study on Kafka, Artemis, Pulsar, and RocketMQ} oleh Md Showkat Hossain Chy, Muhammad Ashfakur Arju, Sri Manjusha Tella, dan Tomas Cerny (2023) \cite{chy2023comparative}. Penelitian ini menyoroti keterbatasan studi evaluasi antrean pesan yang umumnya hanya mengukur kecepatan respons tanpa memperhitungkan biaya konsumsi perangkat keras. Melalui metode pengujian eksperimental terkontrol, penulis menggabungkan pengukuran \textit{throughput} dan latensi dengan instrumentasi pemakaian CPU, alokasi memori, serta aktivitas \textit{garbage collection} pada empat layanan antrean berbasis JVM. Hasil pengujian memperlihatkan bahwa peningkatan \textit{throughput} dan jumlah partisi berbanding lurus dengan lonjakan konsumsi CPU serta memori yang signifikan pada masing-masing \textit{broker}. Walaupun metodologi pengukurannya sangat relevan dalam memadukan metrik kinerja dan profil sumber daya internal, objek penelitian tersebut terbatas pada \textit{broker} distribusi pesan langsung dan belum menyentuh mekanisme eksekusi tugas terjadwal berdasarkan target waktu (\texttt{scheduled\_at}).
	\item \textbf{Unleashing Real-Time Analytics: A Comparative Study of In-Memory Computing vs. Traditional Disk-Based Systems} oleh Semen M. Levin (2024) \cite{levin2024unleashing}. Penelitian ini mengkaji hambatan latensi I/O pada sistem manajemen basis data relasional berbasis disk ketika dihadapkan pada beban pemrosesan waktu nyata. Penulis membandingkan waktu eksekusi dan efisiensi pemrosesan antara penyimpanan \textit{in-memory} menggunakan Redis dan basis data relasional berbasis disk menggunakan PostgreSQL pada berbagai skenario pembacaan, penulisan, serta agregasi data. Hasil eksperimen menunjukkan bahwa penyimpanan data di memori memangkas latensi akses secara drastis dibandingkan eksekusi kueri SQL pada PostgreSQL yang harus melewati lapisan indeks dan penulisan disk. Namun, pengujian Levin berfokus pada beban kueri analitik umum dan belum mengevaluasi perilaku kedua penyimpanan tersebut ketika digunakan sebagai mesin penjadwal tugas konkuren di bawah beban stabil (\textit{uniform load}), lonjakan antrean serentak (\textit{spike load}), maupun skenario penghentian paksa proses aplikasi.
	\item \textbf{Comparative Performance Analysis of RabbitMQ and Kafka Message Queue Systems in Spring Boot and ASP.NET Environments} oleh Filip Kami{\'n}ski, Rados{\l}aw K{\l}onica, dan Beata Pa{\'n}czyk (2025) \cite{kaminski2025comparative}. Penelitian ini meneliti bagaimana variasi beban trafik dan pemilihan lingkungan \textit{runtime} aplikasi memengaruhi kinerja serta pemakaian sumber daya dari sistem antrean pesan yang dijalankan di dalam \textit{container} Docker. Penulis mengukur waktu pemrosesan pesan, \textit{throughput}, serta utilisasi CPU dan RAM pada berbagai skenario pengiriman pesan tunggal maupun masif. Hasil pengujian membuktikan bahwa selain perbedaan mekanisme internal antrean, karakteristik bahasa pemrograman dan \textit{runtime} aplikasi pengelola antrean turut memberikan selisih pemakaian CPU dan memori yang nyata. Temuan ini memperkuat keputusan perancangan pada Tugas Akhir ini untuk mengunci seluruh kandidat strategi penjadwalan di atas satu bahasa pemrograman dan versi \textit{runtime} yang sama (Go v1.26.3) serta batas sumber daya \textit{container} yang identik, sehingga perbedaan performa yang terukur murni mencerminkan efisiensi arsitektur penjadwalannya.
\end{enumerate}

Ringkasan perbandingan kelima penelitian terdahulu tersebut beserta posisi penelitian Tugas Akhir ini disajikan pada Tabel \ref{table:2.literasi}. \par
\begin{longtable}{| b{0.04\textwidth}|p{0.21\textwidth}|p{0.20\textwidth}|p{0.15\textwidth}|p{0.25\textwidth}|}
	\caption{Literasi Penelitian Terdahulu}
	\label{table:2.literasi}\\
	\hline
	\textbf{No.} & \textbf{Judul \& Penulis} & \textbf{Masalah} & \textbf{Metode} & \textbf{Hasil \& Keterbatasan} \\
	\hline
	\endfirsthead
	\hline
	\textbf{No.} & \textbf{Judul \& Penulis} & \textbf{Masalah} & \textbf{Metode} & \textbf{Hasil \& Keterbatasan} \\
	\hline
	\endhead
	1. & \textit{Benchmarking Message Queues} (Maharjan dkk., 2023) \cite{maharjan2023benchmarking} & Pemilihan \textit{message queue} yang tepat untuk beban kerja berbeda antara latensi rendah dan \textit{throughput} tinggi. & Eksperimen komparatif (Redis, RabbitMQ, Artemis, Kafka) dengan OpenMessaging Benchmark. & Redis unggul pada latensi terendah, Kafka unggul pada \textit{throughput}. Belum menguji penjadwalan tertunda (\textit{temporal jitter}) dan konsumsi CPU/RAM. \\ 
	\hline
	2. & \textit{DBOS: A DBMS-Oriented Operating System} (Skiadopoulos dkk., 2022) \cite{skiadopoulos2022dbos} & Kompleksitas pengelolaan status penjadwalan dan pemulihan kegagalan pada sistem terdistribusi. & Eksperimen penjadwalan tugas dan IPC di atas tabel RDBMS berbasis transaksi ACID dan indeks B-Tree. & Penjadwalan berbasis RDBMS memiliki kinerja kompetitif dan mempermudah \textit{fault recovery}. Belum dibandingkan langsung dengan penjadwal \textit{in-memory} dan Redis pada node tunggal. \\ 
	\hline
	3. & \textit{Comparative Evaluation of JVM-Based Message Queue Services} (Chy dkk., 2023) \cite{chy2023comparative} & Evaluasi antrean pesan jarang mengukur konsumsi sumber daya perangkat keras secara bersamaan. & Eksperimen terkontrol pada Kafka, Artemis, Pulsar, dan RocketMQ dengan memantau \textit{throughput}, latensi, CPU, dan RAM. & Peningkatan \textit{throughput} memicu lonjakan pemakaian CPU dan memori yang berbeda tiap arsitektur. Terbatas pada pengiriman pesan langsung berbasis JVM. \\ 
	\hline
	4. & \textit{Unleashing Real-Time Analytics: A Comparative Study of In-Memory Computing vs. Traditional Disk-Based Systems} (Levin, 2024) \cite{levin2024unleashing} & Hambatan latensi I/O pada RDBMS berbasis disk saat memproses beban waktu nyata dibandingkan penyimpanan \textit{in-memory}. & Pengujian komparatif kecepatan baca, tulis, dan agregasi data antara Redis dan PostgreSQL. & Redis memangkas latensi akses secara signifikan dibanding PostgreSQL. Belum menguji mekanisme penjadwalan tugas konkuren pada pola \textit{spike load} dan uji \textit{crash}. \\ 
	\hline
	5. & \textit{Comparative Performance Analysis of RabbitMQ and Kafka in Spring Boot and ASP.NET} (Kami{\'n}ski dkk., 2025) \cite{kaminski2025comparative} & Pengaruh variasi beban kerja dan lingkungan \textit{runtime} terhadap kinerja serta pemakaian CPU/RAM kontainer antrean pesan. & Eksperimen beban tunggal dan masif dalam lingkungan Docker dengan pemantauan waktu respons, CPU, dan memori. & Karakteristik \textit{runtime} aplikasi ikut memengaruhi konsumsi CPU dan RAM selain arsitektur antrean. Hanya menguji RabbitMQ dan Kafka tanpa skenario penjadwalan waktu target. \\ 
	\hline
	6. & \textbf{Penelitian ini} (2026) & Belum adanya perbandingan empiris lintas paradigma (\textit{In-Memory Engine}, PostgreSQL \texttt{SKIP LOCKED}, dan Redis \textit{Sorted Set}) pada \textit{runtime} seragam yang mengukur \textit{jitter}, sumber daya, dan pemulihan \textit{crash}. & Eksperimen kuantitatif komparatif pada \textit{testbed} terkontrol (Go v1.26.3, Docker, Chrony, cAdvisor, Prometheus) dengan skenario \textit{uniform load}, \textit{spike load}, dan uji kegagalan. & Memetakan \textit{trade-off} terukur antara \textit{temporal jitter} ($p50, p95, p99$), \textit{throughput}, konsumsi CPU/RAM/I/O, serta durabilitas pasca-\textit{crash} sebagai dasar rekomendasi pemilihan strategi. \\ 
	\hline
\end{longtable}

Berdasarkan kajian terhadap penelitian-penelitian di atas, terlihat bahwa evaluasi yang ada umumnya hanya membandingkan sesama \textit{message broker} untuk pengiriman pesan langsung \cite{maharjan2023benchmarking, chy2023comparative, kaminski2025comparative} atau membandingkan kecepatan akses Redis dan PostgreSQL pada kueri data umum \cite{levin2024unleashing}. Di sisi lain, pemanfaatan basis data relasional sebagai mesin penjadwal yang tahan gagal \cite{skiadopoulos2022dbos} belum diadu secara langsung dengan antrean memori internal maupun Redis \textit{Sorted Set} di bawah batas komputasi dan \textit{runtime} yang diseragamkan. Celah inilah yang dijawab oleh penelitian Tugas Akhir ini melalui pembangunan \textit{benchmarking testbed} terkontrol untuk menguji \textit{In-Memory Engine}, \textit{Database-Backed Worker} (PostgreSQL), dan \textit{Queue Engine} (Redis) secara berdampingan. \par

\section{Dasar Teori} \label{II.Teori}
Berisi teori/konsep yang berkaitan/digunakan dalam tugas akhir yang dikerjakan. Gunakanlah data melalui buku/jurnal referensi, publikasi tugas akhir, penelitian, buku, dan informasi web yang dapat dipertanggungjawabkan, hindari penggunaan dasar teori melalui tautan Wikipedia, surat kabar, atau portal berita, yang dapat memiliki isi yang tidak bersifat fakta. \par

\subsection{Teori 1} \label{II.Teori1}
Berikut adalah contoh penyisipan tabel menggunakan \verb|\begin{longtable}{}|: \par
	
	\begin{longtable}{|c|c|c|c|}
		\caption{Contoh Tabel}
		\label{table:2.contoh}\\
		\hline
		Col1 & Col2 & Col2 & Col3 \\
		\hline
		\endhead
		1 & 6 & 87837 & 787 \\ 
		\hline
		2 & 7 & 78 & 5415 \\
		\hline
		3 & 545 & 778 & 7507 \\
		\hline
		4 & 545 & 18744 & 7560 \\
		\hline
		5 & 88 & 788 & 6344 \\
		\hline
	\end{longtable}

\subsubsection{Subsubbab} \label{II.Teori1.1}
Berikut adalah contoh subsubbab. Ini adalah level subbab maksimal dalam laporan Tugas Akhir, dan tidak boleh lebih dalam. \par

Gambar \ref{fig:2.contoh} adalah contoh Gambar yang diambil dari internet yang harus dicantumkan sumbernya dan memiliki lisensi Creative Common. Jika gambar adalah milik peneliti lain atau tidak dibuat atau diambil sendiri maka peneliti wajib meminta izin kepada peneliti lain tersebut untuk mencantumkan gambar. Gunakan \verb|\begin{figure}| untuk memasukkan gambar. Gunakan \verb|\caption{[nama caption]}| untuk memberikan caption gambar. Nomor caption akan diurutkan secara otomatis. Jangan lupa untuk melabeli setiap gambar dengan \verb|\label{[nama label]}|, agar bisa direferensi menggunakan \verb|\ref{[nama label]}| \par
\begin{figure}[H] % Kalau menggunakan H, posisi gambar akan tepat dibawah teks
	\centering
	\includegraphics[width=0.6\textwidth]{figure/keyboard.jpg}
	\caption{Contoh gambar dan caption}
	\label{fig:2.contoh}
	{\footnotesize Sumber: Contoh} % Untuk memberikan sumber
\end{figure}

\subsection{Teori 2} \label{II.teori2}
Untuk membuat sebuah rumus persamaan, gunakan kode \verb|\begin{equationcaptioned}| seperti dibawah: \par
	
\begin{equationcaptioned}[eq:2.sederhana]{
	x + 1 = 2
}{
	Rumus sederhana % Caption rumus
}
\end{equationcaptioned}

Teks caption rumus tidak akan muncul di teks, tetapi akan muncul di Daftar Rumus. \par

\begin{equationcaptioned}[eq:2.mae]{
    MAE = \frac{1}{n} \sum_{i=1}^{n} \left| y_i - \hat{y}_i \right|
}{
    Mean Absolute Error (MAE)
}
\end{equationcaptioned}

Berikut adalah contoh penulisan persamaan yang lebih kompleks, yaitu persamaan distribusi normal. \par

\begin{equationcaptioned}[eq:2.mae]{
		P(x) = \frac{1}{{\sigma \sqrt {2\pi } }}e^{{{ - \left( {x - \mu } \right)^2 } \mathord{\left/ {\vphantom {{ - \left( {x - \mu } \right)^2 } {2\sigma ^2 }}} \right. \kern-\nulldelimiterspace} {2\sigma ^2 }}}
	}{
		Distribusi Normal
	}
\end{equationcaptioned}

Jika menuliskan banyak persamaan secara berurutan, gunakan  \verb|\begin{split}|: \par

\begin{equationcaptioned}[eq:2.mae]{
		\begin{split} 
			2x - 5y &=  8 \\ 
			3x + 9y &=  -12
		\end{split}
	}{
		Sistem persamaan linier
	}
\end{equationcaptioned}